\section{Related Work}\label{related}
%\fixme{Raj@Orchid: We need to cite our own works but we will do once we know the outcome from conference.}
This section describes previous works related to multi-hop reasoning (Section \ref{sec:mh}) and fact-checking with KG (Section \ref{sec:fc}). % The multi-hop reasoning section gives an overview of the emergence of reinforcement learning techniques for relation and entity prediction in KGs. The fact-checking section introduces fact-checking works that employ knowledge graphs, explain the fact-checking model or do both.


\subsection{Multi-hop reasoning}\label{sec:mh}
KG reasoning infers new knowledge from existing knowledge in the KG, and detects errors in the KG. Large scale KGs are inherently incomplete. The main effort in KG reasoning research is to ``fill in'' the missing knowledge. This is done by either KG traversal known as multi-hop reasoning \cite{Das2017,  Lv2020, Wan2021} 
%\cite{Das2017, Shen2018, Lin2018, Lv2019, Li2019, Fu2019, Hildebrandt2020, Wan2020, Lei2020, Lv2020, Wan2021, Lao2011, Neelakantan2015, Das2017a, Xiong2017, Chen2018, Wang2019}, 
or by mining logic rules based on KG semantics and using those rules for reasoning \cite{Zhang2019, Qu2021, Minervini2021}.  %{@Gustav, I have removed citations as it was taking too much space due to page limit, hope its okay? Kindly let know this}
%\cite{Galarraga2015, Omran2018, Ho2018, Meilicke2019, Yang2017, Rocktaeschel2017, Sadeghian2019, Qu2021, Minervini2021, Zhang2019, Qu2019}
There are two main directions in multi-hop reasoning: entity prediction and relation prediction. The former looks for a tail entity for a head entity and relation pair
%$$<e_h, r, ?> -> <e_h, r, e_t>$$
and the latter looks for a missing relation between two entities.
%$$<e_h, ?, e_t> -> <e_h, r, e_t>$$
The task of looking for a missing tail entity is more similar to our work of finding a path that leads from a given head entity to the desired (true) tail entity. The major difference is that the proposed approach receives the claimed tail entity as an input as well.

Reasoning in the continuous vector space of KG embeddings was first proposed by Xiong et al. \cite{Xiong2017}. Compared to Das et al. \cite{Das2017a} %\cite{Lao2011, Neelakantan2015, Das2017a} 
which rely on Path Ranking Algorithm (PRA) to produce the initial paths, Xiong et al. \cite{Xiong2017} utilizes a policy-based RL algorithm REINFORCE to learn paths in the KG. Das et al. \cite{Das2017} builds on Xiong et al. \cite{Xiong2017} by directing the task from relation prediction to link prediction, a more complicated task, and also building a model that can handle multiple reasoning entities. %, \cite{Xiong2017} being a model-per-relation task. 
We use both the proposed target entity and relation as input, thus differing from both Das et al. 
\cite{Das2017} as well as Xiong et al. \cite{Xiong2017}.

\subsection{Fact-checking with Knowledge Graph}\label{sec:fc}
%Computational fact-checking generally utilizes machine learning techniques and recently the field has diverted attention towards the interpretability of such machine learning solutions' predictions asking the question "Why?" to combat the machine learning black-box problem.
%For the sake of brevity the overview of fact-checking related work is limited to fact-checking that utilizes knowledge graphs or attempts to explain the fact-checking process or result.
%\subsubsection{Fact-checking with knowledge graphs}
%\textcolor{red}{@Gustav, I have completely removed explainable fact checking below sub section that was without KG and moved it to Introduction, as our work doesnt directly relates with that, kindly read the below section hope its fine..}
Ciampaglia et al. \cite{Ciampaglia2015} framed the problem of fact-checking as finding a connection between the subject and object of a fact statement: assertion is true if there is an edge or short path between the entities in a KG, and false otherwise. Their proposed method makes use of a truth valuation function that accounts for the length of the path between the two entities as well as the generality of intermittent entities, where generality means \textit{how many statements it participates in.}
Shiralkar et al. \cite{Shiralkar2017} builds upon the notion of generality of intermittent entities but considers multiple paths between the subject and object entities as a knowledge stream.
These approaches rely on pre-extracted subgraphs. Inversely, the proposed approach saves resources by integrating relevant subgraph extraction as a part of the explainable classification process.

Previous research is also involved with mining the KG for patterns to be used in classifying a fact statement.
Shi et al. \cite{Shi2016} mined the KG for patterns associated with specific relations. %For example if a city is the capital of a country then very often the parliament will be situated in that city. These patterns are used to classify the fact statement.
Lin et al. \cite{Lin2019} used a supervised pattern discovery algorithm to find patterns relevant to training data. Lin et al. \cite{Lin2018a} builds upon this work by also taking into account the ontological closeness of entities. %, for example: "logician" and "theologician" are ontologically close to "philosopher" and this should be taken into account in pattern mining.
 Unlike these works, this work wholly operates in the continuous vector space and is not formally concerned with patterns, although patterns do arise naturally from the RL aspect.
%Not sure how related Ostrowski is. It uses multi-hop in the title, but is not a classic KG multi-hop task
%\cite{Ostrowski2021} leverages graph format to find connections between relevant sentences to classify a factual statement.
\begin{comment}
\subsubsection{Explainable fact-checking}
\cite{Atanasova2020, Kotonya2020} both formulate the task of explanation generation as an extractive-abstractive summarization task using transformer models introduced in \cite{liu2019}. \cite{Atanasova2020} deals with a political news dataset from \cite{Alhindi2018} and \cite{Kotonya2020} explores facts that require specific expertise with a public health related dataset. Both datasets are annotated with gold standard justifications by journalist. These approaches rely on the annotations which contain the ground truth, unlike our method, which extracts the truth from large scale knowledge graphs.
\cite{Sarrouti2021} evaluates health related claims against scientific articles using transformer models.
%\cite{Atanasova2020, Kotonya2020, Lu2020, Sarrouti2021, Dai2022}
\end{comment}
\par

\noindent \textbf{Explainable fact-checking with Knowledge Graphs: }Ahmadi et al. \cite{Ahmadi2019} used KGs to assess claims and accompany these with human-readable explanations. This is achieved through a combination of logical rule discovery and probabilistic answer set programming (ASP). The KG is mined for rules in the form of triples associated with specific relations and these relations' negations. For fact assessment the heads of these rules are replaced by claim values and the KG is scavenged for corresponding triples. Explanations and an assessment are derived from these triples through ASP. Gad-Elrab et al. \cite{GadElrab2019} used rules to extract explanations from KGs and text resources. Schatz et al. \cite{Schatz2021} mined for patterns for fact-discovery in the biomedical domain to produce human-readable explanations to professionals for drug or treatment discovery. The main difference between these previous approaches and the proposed  approach is that it does not rely on mined patterns but a learned policy and KG embeddings.

